% page 337 du fac-simile (image p-336 du scan)
% bloc 170.2 x 224.7 mm ; marge gauche 31.8 mm
\pgc{10.160mm}{0.000mm}{
\l{\cel{55}329}
\l{}
\l{\sou{lanco}.- I. (\sou{epoki antiqua}).\cel{2}Armo, ek shafto finanta per ferajo}
\l{\cel{3}quadriangula, pintoza, uzata sive quaze flecho quan onu lansas}
\l{\cel{3}manue, sive quaze piquo per qua onu kombatas proxim la adverso.-}
\l{\cel{3}II. (\sou{mezepoko}). Armo kun shafto longa e tre pezoza,finanta}
\l{\cel{3}per ferajo pintoza, per qua la +chevaliero, su-precipitinte}
\l{\cel{3}a sua adverso, probas renversar lu de lua kavalo, e trapikar}
\l{\cel{3}lu. - III. (\sou{nia-epoke}).Armo kun longa shafto lejera, finanta}
\l{\cel{3}per ferajo pintoza, ornita ye banderolo ek stofo, uzata da ula}
\l{\cel{3}korpi di kavalrio (lancieri, dragoni). - DEFIRS.}
\l{}
\l{\sou{lancinar}. (\sou{netrans.}) (\sou{pri doloro pikoza}). Aparar bruske en parto}
\l{\cel{3}di la korpo. - FIS.}
\l{}
\l{\sou{lando}.- Dividuro de regiono geografiala, determinita per limiti}
\l{\cel{3}naturala o politikala. - DE.}
\l{}
\l{\sou{\textquotedbl{}landau\textquotedbl{}}. Veturo quar-rota, kun suspenso-risorti e du tendi,}
\l{\cel{3}forme di berlino, quan onu povas tegar o destegar segun-vole.-}
\l{}
\l{\sou{landgravo}. Titulo di ula suverena princi (olima) di Germania.}
\l{}
\l{\sou{lango}.\cel{2}(\sou{anat.)}\cel{2}Korpo longa, karnoza, moviva, qua esas en la}
\l{\cel{3}kavajo boko. - II. (\sou{teknol.)}\cel{1}Nomo di diversa instrumenti lango-}
\l{\cel{3}forma.FIS.}
\l{}
\l{\sou{langorar}. (\sou{netrans}.)Esar kontinue depresita ed abatita (rispek-}
\l{\cel{3}tive: psikale o korpale). - EFIS.}
\l{}
\l{\sou{langusto}. (\sou{zool.)}\cel{2}Krustaceo analoga a la homardo, ma sen pinci}
\l{\cel{3}e kun longa anteni. - L.\cel{1}\sou{palinurus vulgaris.}\cel{2}- FS.}
\l{}
\l{\sou{lanio.(zool.)}\cel{2}Pasero kun beko kona e hokatra ye sua extremajo,}
\l{\cel{3}qua kombatas uceli plu forta kam su e flugas lansante krii}
\l{\cel{3}akuta. - L.\cel{1}\sou{lanius}. - IL.}
\l{}
\l{\sou{lanjo}. Speco de stofo lana per qua onu envelopas la infanteti}
\l{\cel{3}ankore kelka-monata. - F.}
\l{}
\l{\sou{lansar}.\cel{2}(\sou{trans.}) I. Igar departar tre rapide, aden direciono}
\l{\cel{3}determinita. - II. (\sou{metaf.)}\cel{2}Igar funcionar, igar efikar.-}
\l{\cel{3}DEFIS.}
\l{}
\l{\sou{lantano}. I. (\sou{bot.)}\cel{3}Planto exotika, ek la famiio \textquotedbl{}verbenadei\textquotedbl{},}
\l{\cel{3}kultivata precipue kom apartamento-planto. -L. lantana. -}
\l{\cel{3}II. (\sou{kemio)}\cel{3}Metalo, korpo simpla qua deskompozas aquo e qua}
\l{\cel{3}divenas kombustebla en aero. -\cel{1}\sou{Simb. kem}.: La. - DEFIS.}
\l{}
\l{\sou{lanterno}. Sorto di buxo di qua la parieti esas garnisita ye}
\l{\cel{3}materio plu o min diafana (korno, vitro, papero, e c.) en qua}
\l{\cel{3}onu pozas lumo-fonto shirme de vento. - DEFIS.}
}
